\documentclass[12pt,a4paper]{article}
\usepackage[utf8]{inputenc}
\usepackage[T1]{fontenc}
\usepackage{amsmath,amssymb,amsfonts}
\usepackage{hyperref}
\usepackage{geometry}
\usepackage{booktabs}
\usepackage{amsthm}

\newtheorem{theorem}{Theorem}
\newtheorem{proposition}[theorem]{Proposition}
\newtheorem{lemma}[theorem]{Lemma}
\newtheorem{corollary}[theorem]{Corollary}
\newtheorem{conjecture}[theorem]{Conjecture}
\newtheorem{observation}[theorem]{Observation}
\theoremstyle{remark}
\newtheorem{remark}[theorem]{Remark}

\geometry{a4paper, margin=1in}

\title{Finite-Scale Shadowing: Local $2$-Adic Descent\
and the Two Missing Bridges in an Effective Collatz Bound}
\author{Collatz Crystal Hunter Project}
\date{August 2026}

\begin{document}
\maketitle

\begin{abstract}
We separate two previously conflated mechanisms in an attempted effective
almost-all Collatz bound.  A finite-scale $2$-adic counting argument rigorously
controls forward descent on local scale windows, but iteration requires an open
restart/transport estimate for the endpoint distribution.  Independently, the
tilted reverse model produces the geometric ratio
$\gamma_{\rm rev}=I_{\rm rev}(1.33)/(\log_2 3-1.33)\approx0.993$; attaching this
reverse-growth exponent to the forward exceptional set requires a distinct
reverse-chord encoding hypothesis.  We formulate both missing bridges
explicitly and show that the direct method of types gives a positive local rate
for $Y\asymp N_0^\alpha$ whenever
$1<\alpha<2/\log_2 3\approx1.26186$.  Thus $0.993$ is presently a
reverse-geometric characteristic, not an established exceptional-set exponent.
\end{abstract}

\section*{Units and conventions}
All rate functions, entropies and gains are evaluated consistently in a single
unit; we adopt \emph{bits}. Concretely, the tilt is $2^{s a}$ with $s<1$,
$\Lambda_n(s)=\log_2\rho_n(s)-\log_2\rho_n(0)$ (so $\Lambda(0)=0$, $\Lambda'(0)=2$),
and $I(\sigma)=\sup_s[s\sigma-\Lambda(s)]$ is in bits. Equivalently, in nats
(tilt $e^{sa}$, $s<\ln2$, natural logs) $I_{\mathrm{nat}}(\sigma)=I_{\mathrm{bits}}(\sigma)\ln2$.
The reverse structural gain is $g=\log_2 3-1.33\approx0.25496$ bits/step
$=(\ln3-1.33\ln2)\approx0.17668$ nats/step.  The reverse-geometric ratio
$\gamma_{\rm rev}=I_{\rm rev}/g$ is dimensionless and unit-independent:
$0.2532/0.25498\approx0.993=0.17550/0.17673$.  This unit conversion does not turn
$\gamma_{\rm rev}$ into a forward exceptional-set exponent.  The earlier value
$0.67$ arose from mixing a nats rate with a bits gain and is superseded.  We
normalize $I(2)=0$ throughout.

\section*{Introduction}
Tao's almost-all theorem relies on transporting local information between
scales.  The direct $2$-adic distribution of a finite valuation word is
rigorous, but it remains uniform only while the required cylinder modulus fits
inside the sampled interval.  For forward descent this yields a local window,
not an unrestricted iteration.  The missing ingredient is a restart/transport
lemma controlling the push-forward distribution of surviving endpoints.

A second issue is directional.  The value $1.33<\log_2 3$ describes a reverse
growing chord, whereas a forward segment that remains above $N_0$ satisfies an
upper shift threshold strictly larger than $\log_2 3$.  Therefore the ratio
$I_{\rm rev}(1.33)/(\log_2 3-1.33)$ cannot be assigned to the forward
exceptional set without an additional combinatorial encoding.  This note
records the local rigorous statement and isolates these two open bridges.

\section*{Finite-scale counting}

Let $S_d$ denote the sum of the first $d$ Syracuse valuations.  A valuation
word of total weight $S$ determines one $2$-adic cylinder of relative density
$2^{-S}$ among odd integers, while the number of such words is
$\binom{S-1}{d-1}$.

\begin{proposition}[Finite-scale $2$-adic LDP]
Let $\mathcal I$ contain $M$ odd integers and let $m\ge d$.  Then
\[
 \mathbb P_{N\in\mathcal I}(S_d\le m)
 \le \mathbb P(G_1+\cdots+G_d\le m)
   +O\!\left(\frac{\binom md}{M}\right),
 \qquad G_i\stackrel{\rm iid}{\sim}{\rm Geom}(1/2).
\]
If $m=\sigma d$ and $m\le\log_2M+O(\log d)$, this implies
\[
 \mathbb P_{N\in\mathcal I}(S_d\le\sigma d)
 \le 2^{-dI_2(\sigma)+O(\log d)}.
\]
\end{proposition}
\begin{proof}
Summing the cylinder masses and interval-counting errors over valuation words
gives
\[
 \sum_{S=d}^m\binom{S-1}{d-1}2^{-S}
 +O\!\left(M^{-1}\sum_{S=d}^m\binom{S-1}{d-1}\right).
\]
The first sum is the stated geometric probability and the second word count is
$\binom md$.  The asymptotic bound follows from Cram\'er's estimate and
$I_2(\sigma)+\sigma H_2(1/\sigma)=\sigma$.
\end{proof}

\begin{corollary}[Unconditional local descent window]
Let $\lambda=\log_2 3$, fix $1<\alpha<2/\lambda$, and put
$Y\asymp N_0^\alpha$.  Choose
\[
 \frac{\alpha-1}{2-\lambda}<t\le\frac1\lambda,
 \qquad d=\lfloor t\log_2N_0\rfloor,
 \qquad \sigma=\lambda+\frac{\alpha-1}{t}<2.
\]
The proportion of odd $N\asymp Y$ whose first $d$ accelerated odd iterates all
exceed $N_0$ is
\[
 \le N_0^{-tI_2(\sigma)+o(1)}.
\]
The rate is positive for every $\alpha<2/\log_2 3\approx1.26186$ and
degenerates at the endpoint.
\end{corollary}
\begin{proof}
If $x_k>N_0$ for $0\le k\le d$, then the exact identity
\[
 x_d=N3^d2^{-S_d}\prod_{k<d}\left(1+\frac1{3x_k}\right)
\]
gives
\[
 S_d<d\lambda+\log_2(N/N_0)+O(d/N_0).
\]
For $N\asymp N_0^\alpha$ this threshold is $\sigma d+o(d)$.  The condition
$t\lambda\le1$ makes the required cylinder modulus fit inside the sampled
interval, and the lower bound on $t$ ensures $\sigma<2$.
\end{proof}

\section*{The two open bridges}

\begin{lemma}[Within-word endpoint transport]
Fix an admissible valuation word of total weight $S$ and length $d$.  There are
odd integers $\rho_w,y_w$ such that
\[
 N=\rho_w+2^{S+1}q,
 \qquad
 \operatorname{Syr}^d(N)=y_w+2\cdot3^dq.
\]
The barrier conditions on all intermediate iterates restrict $q$ to an integer
interval.  Hence in every odd residue class modulo $2^m$ the endpoint count
differs from the uniform value by at most one.
\end{lemma}
\begin{proof}
Apply the affine word identity at every prefix.  Exact valuation among odd
starts fixes a class modulo $2^{S+1}$, and all prefix maps are increasing in
$q$.  Since $3^d$ is invertible modulo $2^{m-1}$, the endpoint progression
permutes odd residue classes modulo $2^m$.
\end{proof}

\begin{proposition}[Exact boundary-layer spectrum]
Let $K=2^M$, use the odd coordinate $N=2z+1$, and set
$\mathcal B_{d,M}=\{z:S_d(2z+1)\le M\}$.  For an exact valuation word $w$
of weight $S$, let $r_w\pmod{2^S}$ be its residue in the $z$ coordinate.
Then, for $h\ne0$ and $v=v_2(h)$,
\[
 \frac{\widehat{\mathbf1_{\mathcal B_{d,M}}}(h)}K
 =\sum_{S=\max\{d,M-v\}}^M2^{-S}
   \sum_{\substack{w\in\mathbb N^d\\|w|=S}}
   \exp(-2\pi i h r_w/K).
\]
In particular, odd frequencies see only the top layer $|w|=M$.
\end{proposition}
\begin{proof}
The exact word condition is
$3^dN+c_w\equiv2^S\pmod{2^{S+1}}$, not merely divisibility by $2^S$.
Thus $z\equiv r_w\pmod{2^S}$.  The Fourier transform of this cylinder
vanishes unless $2^{M-S}\mid h$ and otherwise has normalized value
$2^{-S}\exp(-2\pi i h r_w/K)$.  Sum over the disjoint words with $S\le M$.
\end{proof}
\begin{lemma}[Top-layer offset injectivity]
Fix $d\le M$ and distinct words $w,w'\in\mathbb N^d$ of total weight $M$.
Then
\[
 r_w\not\equiv r_{w'}\pmod{2^{M-1}}.
\]
In particular, $A_{\mathrm{ant}}=0$ on every fixed $(d,M)$ top layer.
This statement does not compare words of different lengths.
\end{lemma}
\begin{proof}
If $S_j$ and $S'_j$ are the partial sums, then
\[
 c_w=\sum_{j=0}^{d-1}3^{d-1-j}2^{S_j}.
\]
At the smallest exponent in the symmetric difference of the two partial-sum
sets, the difference $c_w-c_{w'}$ has an odd leading coefficient; hence
$v_2(c_w-c_{w'})<M$.  The CRT formula
$\rho_w\equiv(2^M-c_w)3^{-d}\pmod{2^{M+1}}$ then shows that equality of
$r_w$ and $r_{w'}$ modulo $2^{M-1}$ is impossible.
\end{proof}

\begin{lemma}[Spectral Transversality of Resonant Sets; falsified as a decay mechanism]
\label{lem:spectral_transversality}
Let $R_{\text{end}} = \{h : \|h 3^{d_{\text{end}}} / 2^M\| \le A/2^M\}$ be the set of resonant frequencies for the endpoint distribution, and $R_{\text{bad}} = \{h : \|h 3^{d_{\text{bad}}} / 2^M\| \le A/2^M\}$ be the resonant set for the bad-set boundary layer.
While the arithmetic fact $d_{\text{end}} \not\equiv d_{\text{bad}} \pmod{2^{M-2}}$ holds, enforcing that the central peaks do not intersect, this separation does not imply smallness of the total error $\Delta_b$. The high-frequency noise cross-terms generate macroscopic peaks of magnitude $\sim 0.43$ independent of the resonance alignment, exceeding the required $2^{-B/2}$ threshold by a factor of 885.
\end{lemma}

\begin{remark}[Falsification of intra-layer anomalous cancellation]
Extensive exact numerical DP tracking over the bucket constraints (up to $M \le 10^4$) explicitly falsifies any hypothesis of anomalous intra-layer Fourier cancellation. Within any fixed, macroscopic slice of length $d \approx M/2$, the Fourier coefficients uniformly obey $\theta \approx 0.47 \approx 0.5$, mirroring an unstructured random walk. Thus, the decay in $\Delta_b$ is not driven by the endpoint distribution becoming artificially smooth, but rather by its sharp peaks systematically avoiding the structural resonances of the bad set.
\end{remark}

\begin{remark}[Finite diagnostic and why pointwise transport is too strong]
At the full next-event modulus, many endpoint residue classes are empty, so
pointwise Haar domination is numerically false and unnecessary.  For
$\alpha=1.10$ and $B=16,18,20$, exhaustive enumeration gives mass-weighted
bad-event errors
\[
 0.003506,\qquad0.001901,\qquad0.000887,
\]
while multiplication by $2^{B/2}$ gives $0.898,0.973,0.908$.  At $B=20$ the
weighted absolute Fourier majorant is $0.05268$, versus a signed aggregate
cancellation ratio $0.00911$.  Three scales identify the target mechanism but
do not prove decay.  Separate exhaustive scale-kernel computations give
bucket-weighted discrepancies of order $0.5\%$--$1.6\%$.
\end{remark}

\begin{remark}[Reverse-chord encoding: superseded]
\label{rem:reverse_encoding_superseded}
The Reverse-chord Encoding Hypothesis was previously formulated as a
potential bridge to attach $\gamma_{\rm rev} \approx 0.993$ to the forward
exceptional set. Numerical experiments (E9) revealed a fundamental
information-theoretic barrier: the entropy of reverse chords with mean
shift $4/3$ is $h(4/3) \approx 1.0817$ bits/step, while the target space
requires $\log_2 3 \approx 1.585$ bits/step. The resulting Hausdorff
dimension deficit ($D_{\rm rev} \approx 0.811 < D_{\rm fwd} \approx 0.950$)
makes 100\% coverage algebraically impossible. This route is closed.
\end{remark}

\begin{proposition}[Conditional reverse-geometric ratio]
The reverse thermodynamic model assigns to a mean-$1.33$ growing chord the
cost-per-gain ratio
\[
 \gamma_{\rm rev}
 =\frac{I_{\rm rev}(1.33)}{\log_2 3-1.33}
 \approx\frac{0.2532}{0.25498}\approx0.993.
\]
This is a characteristic of reverse growth geometry.  It becomes a candidate
forward exceptional-set exponent only under both reverse LDP shadowing and the
Reverse-chord Encoding Conjecture.
\end{proposition}

\begin{remark}[Why one shadowing statement is insufficient]
Forward non-descent gives an upper threshold
\[
 \frac{S_d}{d}<\log_2 3+\frac{\log_2(N/N_0)}d+O(N_0^{-1}) \qquad \text{(which is } >\log_2 3\text{)},
\]
whereas the value $1.33<\log_2 3$ belongs to reverse growth.  Thus reverse LDP
shadowing alone does not imply a forward exceptional-set estimate.  The
restart/transport and reverse-encoding problems are logically distinct.
\end{remark}

\begin{remark}[Honest status]
The finite-scale counting proposition and local descent corollary are
unconditional.  No global power-law exceptional-set bound is proved here.
The previously quoted residual densities at $N_0=2^{68}$ are only scale
illustrations under the stronger, unproved reverse-encoding hypothesis.  A
density statement would still not imply the pointwise Collatz conjecture.
\end{remark}

\end{document}
